% !TeX root = ../../Book3.tex
% 纲要标题：## 第14章　Kähler 微分与无穷小结构
% 纲要位置：工作记录/计划纲要.md，第 2400 行。
% 纲要细目（写作范围）：
% > 本章研究普遍导子及环映射的无穷小结构；带指定导子的微分环、微分理想与微分方程另见附录 E；两种主题不混称。

\chapter{Kähler 微分与无穷小结构}
\label{b3:ch:14}
\BookChapterNumber{b3:ch:14}

\begin{chapterabstract}
本章构造 Kähler 微分模与普遍导子，证明商环、复合映射、局部化及基变换公式，并用 Jacobian 计算切空间。微分还检测域扩张的可分性和平方零提升的唯一性；平展与光滑分别需要进一步的平坦性及提升存在条件。
\end{chapterabstract}

\begin{learninggoals}
\begin{itemize}
\item 从生成元与关系构造微分模，并把导子解释为平方零扩张中的截面。
\item 正确使用两条微分正合列，保留左端未必单射这一限制。
\item 由 Jacobian 矩阵计算有理点切空间，区分局部环余切空间与相对微分纤维。
\item 证明有限域扩张的微分消失判据，计算纯不可分扩张的微分。
\item 分清提升的存在性、唯一性与平坦性，并在标准例子中实际构造提升。
\end{itemize}
\end{learninggoals}

设 \(k\) 为域。抛物线的坐标环 \(B=k[x,y]/(y-x^2)\) 也可写成 \(k[t]\)，其中 \(x\) 对应 \(t\)，\(y\) 对应 \(t^2\)。对平面方程作形式微分，得到 \(\mathrm{d}y=2x\cdot\mathrm{d}x\)；换成后一种展示，则只需以 \(\mathrm{d}t\) 为基。这两种计算应当描述同一个对象。因此，对选定方程求导还不足以作为定义：我们需要由结构同态 \(k\to B\) 本身决定的微分模，使所有函数的一阶变化及其关系都能在其中表达，而不依赖选用了哪些生成元和方程。

\ProseParagraphBreak

一阶变化还可以用平方为零的无穷小量测试：把 \(x\) 送到 \(x+\varepsilon a\)，忽略 \(\varepsilon^2\) 后，乘法相容性恰好给出导子法则。于是微分不仅计算切方向，也能探测一个映射是否允许唯一提升、是否存在提升。比较这些条件时，平坦性和有限展示又承担不同作用，不能仅凭 Jacobian 的外观互相替代。

\ProseParagraphBreak

预备知识是多项式商、模与张量积、局部化及有限域扩张。\chapref{b3:ch:12}的平坦性判据用于比较非分歧与平展，\chapref{b3:ch:07}的坐标环和双数例子则提供几何背景。光滑性所需的一阶定义和计算在本章给出。

\ProseParagraphBreak
微分模及两条正合列可对照松村英之的《Commutative Ring Theory》\cite[\S25, pp. 191--195]{Matsumura1989CommutativeRingTheory}。具体多项式商的微分可以直接计算，再经基变换和局部化核对；普遍性质解释了这些计算为何相容。


% !TeX root = ../../Book3.tex
% 纲要标题：### 14.1 导子与微分模
% 纲要位置：工作记录/计划纲要.md，第 2404 行。
% 纲要细目（写作范围）：
% * 基环上线性导子
% * 微分模的泛性质
% * 多项式环的微分模

\section{导子与微分模}
\label{b3:sec:1401}


把方程中的变量作一个平方为零的改动，乘法展开时只剩常数项与一次项。导子记录的正是这个一次变化。微分模把所有可能的导子合并成一个普遍对象，使后面的商环、局部化和切空间计算都能从同一构造出发。这里研究的是代数本身允许的导子，不是在环上预先指定某一个微分算子。

\subsection{基环上的线性导子}
\label{b3:subsec:1401:01}

\begin{definition}\label{b3:def:relative-derivation}
设 \(\varphi:A\rightarrow B\) 为交换环同态，\(M\) 为 \(B\) 模。如果映射 \(D:B\rightarrow M\) 对所有 \(b,c\in B\)、\(a\in A\) 都满足
\[
D(b+c)=D(b)+D(c),\qquad D(bc)=b\cdot D(c)+c\cdot D(b),\qquad D\bigl(\varphi(a)\bigr)=0,
\]
就称 \(D\) 为 \(A\) 上的\term[daozi]{导子}[derivation]。全体这样的映射记为 \(\operatorname{Der}_A(B,M)\)。点态相加及 \((b\cdot D)(c)=b\cdot D(c)\) 使它成为 \(B\) 模。
\symidx{\(\operatorname{Der}_A(B,M)\)：从代数到模的相对导子}
\end{definition}

Leibniz 法则给出 \(D(1)=2D(1)\)，故 \(D(1)=0\)，并由归纳得到 \(D(b^n)=nb^{n-1}D(b)\)。又
\(D\bigl(\varphi(a)\cdot b\bigr)=\varphi(a)\cdot D(b)\)，所以导子确为 \(A\) 线性映射；它通常不是 \(B\) 线性的，否则 \(D(b)=b\cdot D(1)=0\)。例如 \(\mathrm{d}/\mathrm{d}x:k[x]\rightarrow k[x]\) 是 \(k\) 导子，而 \((\mathrm{d}/\mathrm{d}x)(x)=1\)，并不满足 \(k[x]\) 线性。

指定基环，就是指定哪些元素必须视为常数。设 \(k\) 为域，同一个环 \(B=k(t)\) 相对于 \(k\) 可以有非零导子：按商法则令
\[
D(f/g)=\frac{f'g-fg'}{g^2}\qquad(f,g\in k[t],\ g\ne0),
\]
就得到 \(D(t)=1\) 的 \(k\) 导子。若改为相对于 \(k(t)\) 本身，则 \(B\) 的每个元素都必须被零化，导子只能为零。

\ProseParagraphBreak
若 \(\operatorname{char}B=p>0\)，则任何导子都满足 \(D(b^p)=0\)。特别地，对特征 \(p\) 的域 \(k\)，在 \(k[x]\) 上仍有 \(D(x)=1\) 的形式求导，但所有 \(k\) 导子都零化 \(x^p\)，尽管 \(x^p\notin k\)。因此被导子零化并不意味着来自基域；导子无法检测 \(p\) 次幂方向，这也预示了不可分扩张中的特殊现象。

对交换环 \(B\)，平方零的变化可以先在双数环 \(B[\varepsilon]/(\varepsilon^2)\) 中计算。每个元素唯一写成 \(b+m\varepsilon\)，而
\[
(b+m\varepsilon)(c+n\varepsilon)=bc+(bn+cm)\varepsilon.
\]
对映射 \(D:B\to B\)，\(b\mapsto b+D(b)\varepsilon\) 保持乘法，恰好要求 \(D\) 满足 Leibniz 法则。若变化量取值于一般 \(B\) 模 \(M\)，同样令两项变化量的乘积为零，就得到下面的构造。

\begin{proposition}\label{b3:prop:derivation-square-zero}
设 \(\varphi:A\to B\) 为交换环同态，\(M\) 为 \(B\) 模，\(D:B\to M\) 为映射。在加法群 \(B\oplus M\) 上定义
\[
(b,m)(c,n)=(bc,bn+cm).
\]
这给出一个单位为 \((1,0)\) 的交换环，并通过结构同态
\[
A\longrightarrow B\oplus M,\qquad a\longmapsto\bigl(\varphi(a),0\bigr)
\]
成为 \(A\) 代数。理想 \(0\oplus M\) 的平方为零。映射 \(b\mapsto\bigl(b,D(b)\bigr)\) 是投影 \(B\oplus M\to B\) 的 \(A\) 代数截面，当且仅当 \(D\in\operatorname{Der}_A(B,M)\)。
\end{proposition}
\begin{proof}
三因子乘积的第二坐标都是 \(bc\,u+b d\,n+c d\,m\)，其中第三因子为 \((d,u)\)，故结合律成立；交换律、分配律和单位可直接由公式检验。所列结构映射是环同态，投影到 \(B\) 也保持这一 \(A\) 代数结构。平方为零来自两因子的第一坐标同时为零。截面必须有 \(\bigl(b,D(b)\bigr)\) 的形式，保持加法、乘法以及 \(A\) 的像分别就是导子的三个条件。
\end{proof}

这个环称为平方零扩张 \(B\oplus M\)。它把“无穷小变化”变成确切的环运算；这里的无穷小并非一个任意小的实数，而是乘积被规定为零的理想元素。

\subsection{微分模的泛性质}
\label{b3:subsec:1401:02}

目标模 \(M\) 改变时，\(\operatorname{Der}_A(B,M)\) 也随之改变。我们希望先用一个固定的 \(B\) 模记录导子的关系，再通过从这个模到 \(M\) 的线性映射，给每种一阶变化指定实际的值。

\begin{theorem}\label{b3:thm:universal-differentials}
设 \(A\) 为交换环，\(B\) 为交换 \(A\) 代数。存在 \(B\) 模 \(\Omega_{B/A}\) 与 \(A\) 导子 \(\mathrm{d}:B\rightarrow\Omega_{B/A}\)，使对任意 \(B\) 模 \(M\)，映射
\[
\operatorname{Hom}_B(\Omega_{B/A},M)\longrightarrow
\operatorname{Der}_A(B,M),\qquad h\longmapsto h\circ \mathrm{d}
\]
为自然双射。二者组成的对象在唯一的保持 \(\mathrm{d}\) 的同构意义下唯一，称为 \term[kahler-weifen-mo]{Kähler 微分模}[module of Kähler differentials]；\(\mathrm{d}\) 称为\term[pubian-daozi]{普遍导子}[universal derivation]。
因此，对任意 \(D\in\operatorname{Der}_A(B,M)\)，都有下列交换图；\(\mathrm{d}\) 和 \(D\) 是 \(A\) 导子，虚线 \(h\) 是唯一的 \(B\) 线性映射：
\[
\begin{tikzcd}[column sep=large,row sep=large]
B\arrow[r,"\mathrm{d}"]\arrow[dr,"D"']
&\Omega_{B/A}\arrow[d,dashed,"h"]\\
&M
\end{tikzcd}
\]
\symidx{\(\Omega_{B/A}\)、\(\mathrm{d}_{B/A}\)：Kähler 微分模与普遍导子}
\end{theorem}
\begin{proof}
记结构同态为 \(\varphi:A\to B\)。取以集合 \(B\) 为指标的自由 \(B\) 模 \(F=\bigoplus_{b\in B}Be_b\)，模掉下列元素生成的子模：
\[
e_{b+c}-e_b-e_c,\qquad
 e_{bc}-be_c-ce_b,\qquad e_{\varphi(a)}
\quad(b,c\in B,\ a\in A).
\]
记所得商为 \(\Omega_{B/A}\)，令 \(\mathrm{d}b\) 为 \(e_b\) 的像。三个关系使 \(\mathrm{d}\) 满足导子的三个公理。对任意导子 \(D\)，从自由模出发的唯一线性映射 \(e_b\mapsto D(b)\) 恰好零化这些关系，故下降为 \(h:\Omega_{B/A}\rightarrow M\)。因各 \(\mathrm{d}b\) 生成微分模，这个 \(h\) 唯一。

若另一对 \((\Omega',\mathrm{d}')\) 也有此性质，分别应用泛性质得到 \(u:\Omega\rightarrow\Omega'\)、\(v:\Omega'\rightarrow\Omega\)，满足 \(u\mathrm{d}=\mathrm{d}'\)、\(v\mathrm{d}'=\mathrm{d}\)。于是 \(vu\mathrm{d}=\mathrm{d}\)，唯一性给出 \(vu=1\)，同理 \(uv=1\)。对 \(M\) 的同态作复合即说明双射的自然性。
\end{proof}

构造中的 \(\mathrm{d}b\) 先作为形式符号记录变化，并只施加导子公理要求的关系；选定 \(B\) 线性映射 \(h:\Omega_{B/A}\to M\) 后，这些符号才取得值 \(h(\mathrm{d}b)=D(b)\)。若 \(B\) 由 \(b_i\) 作为 \(A\) 代数生成，则 \(\Omega_{B/A}\) 由 \(\mathrm{d}b_i\) 作为 \(B\) 模生成，因为乘法法则能把任意多项式的微分化成这些元素的线性组合。生成元之间的代数关系则会带来微分关系。

\subsection{多项式环的微分模}
\label{b3:subsec:1401:03}

\begin{proposition}\label{b3:prop:polynomial-differentials}
设 \(A\) 为交换环，\(n\ge0\) 为整数，\(P=A[X_1,\ldots,X_n]\)。则
\[
\Omega_{P/A}=\bigoplus_{i=1}^nP\,\mathrm{d}X_i,
\qquad \mathrm{d}f=\sum_{i=1}^n\dfrac{\partial f}{\partial X_i}\,\mathrm{d}X_i.
\]
对任意 \(P\) 模 \(M\) 和元素 \(m_1,\ldots,m_n\in M\)，存在唯一 \(A\) 导子 \(D:P\rightarrow M\) 满足 \(D(X_i)=m_i\)。
\end{proposition}
\begin{proof}
在单项式上反复使用 Leibniz 法则，任何导子都必须满足所列公式。反向用该公式定义 \(D(f)=\sum_i(\partial f/\partial X_i)m_i\)，多项式形式偏导的乘法公式直接给出 Leibniz 法则，且它零化 \(A\)，所以存在性也成立。于是自由模 \(\bigoplus_iP\,\mathrm{d}X_i\) 具有普遍导子的泛性质，等同于 \(\Omega_{P/A}\)。也可用导子 \(\partial/\partial X_j\) 检查线性无关：它将假定关系 \(\sum_if_i\,\mathrm{d}X_i=0\) 送到 \(f_j=0\)。
\end{proof}

例如在 \(k[x,y]\) 中，\(\mathrm{d}(x^2y)=2xy\,\mathrm{d}x+x^2\,\mathrm{d}y\)。若再施加关系 \(y^2=x^3\)，则必有 \(2y\,\mathrm{d}y-3x^2\,\mathrm{d}x=0\)；下一节将证明怎样获得全部这样的关系。对形式幂级数环则需要额外区分连续导子与任意代数导子，本章的多项式结论不未经证明地推广到无限级数。

% !TeX root = ../../Book3.tex
% 纲要标题：### 14.2 微分模的正合列与基变换
% 纲要位置：工作记录/计划纲要.md，第 2410 行。
% 纲要细目（写作范围）：
% * 商环的余法正合列
% * 方程微分给出的展示；尖点微分模的非零挠元及其零化理想
% * 复合环映射的微分序列
% * 局部化及基变换

\section{微分模的正合列与基变换}
\label{b3:sec:1402}


给定代数的生成元以后，导子的值不能任意选择：它必须保持方程的微分关系。两条基本正合列分别处理“添加关系”和“改变基环”，局部化与基变换公式则使这些计算能够转到所需的局部环或基域。

\subsection{商环的余法正合列}
\label{b3:subsec:1402:01}

商映射 \(B\to B/J\) 把 \(f\in J\) 宣布为零，因而微分中也必须添加 \(\mathrm{d}f=0\) 的关系。两条方程乘积的微分仍含方程因子，转到 \(B/J\) 后就为零，所以 \(J^2\) 中的元素不会提供新的这一阶关系；记录这些方程类的是 \(J/J^2\)。

\begin{theorem}[余法正合列]\label{b3:thm:conormal-sequence}
设 \(A\) 为交换环，\(B\) 为交换 \(A\) 代数、\(J\subseteq B\) 为理想、\(C=B/J\)。有 \(C\) 模正合列
\[
J/J^2\xrightarrow{\delta}C\otimes_B\Omega_{B/A}
\longrightarrow\Omega_{C/A}\longrightarrow0,
\qquad \delta(f+J^2)=1\otimes \mathrm{d}f.
\]
模 \(J/J^2\) 称为相应商映射的\term[yufa-mo]{余法模}[conormal module]。映射 \(\delta\) 一般不单射。
\end{theorem}
\begin{proof}
若 \(f\in J^2\)，将它写成有限和 \(\sum_i u_iv_i\)、\(u_i,v_i\in J\)，则 \(\mathrm{d}f=\sum_i(u_i\,\mathrm{d}v_i+v_i\,\mathrm{d}u_i)\) 在张量到 \(C\) 后为零，故 \(\delta\) 良定义。又 \(\mathrm{d}(bf)=b\,\mathrm{d}f+f\,\mathrm{d}b\)，第二项在 \(C\) 上为零，所以 \(\delta\) 是 \(C\) 线性的。

将 \(C\otimes_B\Omega_{B/A}\) 模掉全部 \(1\otimes \mathrm{d}f\)、\(f\in J\)，记商为 \(Q\)。映射 \(b+J\mapsto[1\otimes \mathrm{d}b]\) 良定义，因为改变代表元只增加一个被模掉的微分。它是 \(A\) 导子。若 \(D:C\rightarrow M\) 是任意 \(A\) 导子，复合 \(B\rightarrow C\) 后的导子对应一个 \(C\otimes_B\Omega_{B/A}\rightarrow M\)，且零化各 \(\mathrm{d}f\)，故唯一通过 \(Q\)。于是 \(Q\) 有 \(\Omega_{C/A}\) 的泛性质，得到所需正合列。
\end{proof}

\begin{corollary}\label{b3:cor:differential-presentation}
设 \(A\) 为交换环，\(f_1,\ldots,f_s\in A[X_1,\ldots,X_n]\)，\(B=A[X_1,\ldots,X_n]/(f_1,\ldots,f_s)\)，\(x_i\) 为变量的像。则有展示
\[
B^s\xrightarrow{J_f^{\mathsf t}} B^n\longrightarrow\Omega_{B/A}\longrightarrow0,
\qquad (J_f)_{ij}=\dfrac{\partial f_i}{\partial X_j}(x).
\]
因此有限展示交换代数的微分模为有限展示模。
特别地，若 \(k\) 为特征不等于 \(2,3\) 的域，取 \(B=k[x,y]/(y^2-x^3)\)，并以 \(x,y\) 兼记变量在商中的像，则
\[
\omega=2x\,\mathrm{d}y-3y\,\mathrm{d}x
\]
是 \(\Omega_{B/k}\) 中的非零挠元，且 \(y\omega=x^2\omega=0\)。
\end{corollary}
\begin{proof}
对 \(P=A[X]\) 的商应用余法正合列，\(\Omega_{P/A}\) 为以 \(\mathrm{d}X_i\) 为基的自由模。若 \(f=\sum_i g_if_i\)，在 \(\Omega_{P/A}\) 张量到 \(B\) 后有 \(\mathrm{d}f=\sum_i g_i\,\mathrm{d}f_i\)，因为含 \(f_i\,\mathrm{d}g_i\) 的项消失。因此全部微分关系由各 \(\mathrm{d}f_i\) 生成。映射 \(B^s\to B^n\) 把第 \(i\) 个标准基向量送到 \(\mathrm{d}f_i\) 的坐标列 \(\bigl(\frac{\partial f_i}{\partial X_1}(x),\ldots,\frac{\partial f_i}{\partial X_n}(x)\bigr)^{\mathsf t}\)。通常 Jacobian 矩阵将第 \(i\) 个方程的偏导放在第 \(i\) 行，这里将它们放在第 \(i\) 列，所以所需矩阵是其转置。

在尖点的情形中，上述展示给出
\[
\Omega_{B/k}=(B\,\mathrm{d}x\oplus B\,\mathrm{d}y)/B(2y\,\mathrm{d}y-3x^2\,\mathrm{d}x).
\]
由 \(y^2=x^3\)，有
\[
y\omega=x(2y\,\mathrm{d}y-3x^2\,\mathrm{d}x)=0,
\qquad x^2\omega=y(2y\,\mathrm{d}y-3x^2\,\mathrm{d}x)=0.
\]
为证 \(\omega\ne0\)，由\cref{b3:prop:cusp-normalization}将 \(B\) 识别为 \(k[t^2,t^3]\subseteq k[t]\)，其中 \(x=t^2\)、\(y=t^3\)。若 \(\omega=0\)，它在自由模中的代表必须是关系生成元的一个 \(B\) 倍数，因此存在 \(h\in B\) 满足
\[
-3y=-3hx^2,\qquad 2x=2hy.
\]
由于 \(2\) 在 \(k\) 中可逆，后一等式在分式域 \(k(t)\) 中迫使 \(h=t^{-1}\)，但 \(B\subseteq k[t]\) 不含此元素，矛盾。因此 \(\omega\) 非零，而 \(y\ne0\) 是整环 \(B\) 的非零元素，所以上面的零化等式也证明它是挠元。
\end{proof}

余法映射不单射，表示非零的方程类也可能没有留下非零的微分关系。例如在特征 \(p\) 的域上取 \(B=k[t]\)、\(J=(t^p)\)。由于 \(t^p\notin(t^{2p})\)，它在 \(J/J^2\) 中的类非零，却满足 \(\delta(t^p)=p t^{p-1}\mathrm{d}t=0\)。这一类生成整个余法模，所以此例的余法映射甚至为零。

\subsection{复合环映射的微分序列}
\label{b3:subsec:1402:02}

对 \(A\to B\to C\)，从相对于 \(A\) 求导改为相对于 \(B\) 求导，意味着将 \(B\) 的像也视为常数。因此要在 \(\Omega_{C/A}\) 中进一步模掉这些元素的微分，末端微分模就成为中间微分模的商。

\begin{theorem}\label{b3:thm:differential-transitivity}
设 \(A\rightarrow B\xrightarrow{g}C\) 为交换环同态，则有 \(C\) 模正合列
\[
C\otimes_B\Omega_{B/A}\xrightarrow{\alpha}\Omega_{C/A}
\xrightarrow{\beta}\Omega_{C/B}\longrightarrow0,
\]
其中 \(\alpha(c\otimes \mathrm{d}b)=c\,\mathrm{d}\bigl(g(b)\bigr)\)、\(\beta(\mathrm{d}_{C/A}c)=\mathrm{d}_{C/B}c\)。
\end{theorem}
\begin{proof}
复合 \(B\rightarrow C\xrightarrow{\mathrm{d}_{C/A}}\Omega_{C/A}\) 是 \(A\) 导子，故给出 \(\alpha\)。在 \(\Omega_{C/A}\) 中模掉各 \(\mathrm{d}\bigl(g(b)\bigr)\) 生成的子模后，普遍导子就零化 \(B\) 的像，因此成为 \(B\) 导子。任意 \(B\) 导子首先是 \(A\) 导子，并零化这些指定元素，所以唯一经过此商。该商即 \(\Omega_{C/B}\)，证明正合性。
\end{proof}

首箭头可能不单射。例如 \(A=k\)、\(B=k[t^p]\subseteq C=k[t]\)、\(\operatorname{char}k=p\) 时，\(t^p\) 是 \(B\) 的一个多项式生成元，所以 \(\mathrm{d}_{B/k}(t^p)\ne0\)；但它映到 \(\Omega_{C/k}\) 中时变成 \(\mathrm{d}_{C/k}(t^p)=0\)。

\subsection{微分模的局部化与基变换}
\label{b3:subsec:1402:03}

只在上环 \(B\) 中倒置乘法集 \(S\) 时，基环仍为 \(A\)。分母 \(s\in S\) 的微分一般不为零，所以分式的微分需要包含 \(\mathrm{d}s\) 项。若分母来自基环中的乘法集 \(T\)，它们本来就必须被零化；同步倒置基环与上环后，这些分母的逆也成为常数，此时相对基环是 \(T^{-1}A\)。

\begin{proposition}\label{b3:prop:differential-localization}
设 \(\varphi:A\to B\) 为交换环同态，\(S\subseteq B\) 为乘法集。则有 \(S^{-1}B\) 模自然同构
\[
\Omega_{S^{-1}B/A}\cong S^{-1}\Omega_{B/A},
\qquad \mathrm{d}(b/s)=\dfrac{\mathrm{d}b}{s}-\dfrac{b\,\mathrm{d}s}{s^2}.
\]
其中 \(b\in B\)、\(s\in S\)。
若 \(T\subseteq A\) 为乘法集，记 \(T^{-1}B\) 为 \(B\) 对 \(\varphi(T)\) 的局部化，则还有 \(T^{-1}B\) 模自然同构
\[
\Omega_{T^{-1}B/T^{-1}A}\cong T^{-1}\Omega_{B/A}.
\]
右侧按 \(A\) 在 \(\Omega_{B/A}\) 上的作用局部化。
\end{proposition}
\begin{proof}
对取值于任意 \(S^{-1}B\) 模 \(M\) 的 \(A\) 导子 \(D:B\rightarrow M\)，若它能延拓为 \(\widetilde D:S^{-1}B\to M\)，则等式 \(\widetilde D(s\cdot s^{-1})=0\) 迫使其满足上述分式公式。按此公式定义延拓，检查良定义：若 \(b/s=b'/s'\)，则存在 \(t\in S\) 使 \(t(s'b-sb')=0\)。取导数并在 \(M\) 中将 \(t,s,s'\) 视为单位；由于 \(s'b-sb'=0\) 在局部化后成立，得到
\[
s'Db+bDs'-sDb'-b'Ds=0.
\]
连同 \(s'b=sb'\)，整理即得两个分式公式相等。对分式加法和乘法展开可得加法性与 Leibniz 法则：乘法中的导数项恰分成对第一个和第二个因子的两部分。故延拓存在且唯一，由微分模的泛性质得到第一个同构。

对于同步局部化，取任意 \(T^{-1}B\) 模 \(M\)。每个 \(A\) 导子 \(D:B\to M\) 同样唯一延拓为 \(\widetilde D:T^{-1}B\to M\)。对 \(t\in T\)，有
\[
\widetilde D\bigl(\varphi(t)^{-1}\bigr)
=-\varphi(t)^{-2}\cdot D\bigl(\varphi(t)\bigr)=0.
\]
它又零化 \(\varphi(A)\)，所以零化 \(T^{-1}A\) 的像，因而是 \(T^{-1}A\) 导子。反过来，限制一个 \(T^{-1}A\) 导子便得到 \(A\) 导子。两种操作互逆，再用微分模的泛性质得到第二个同构。
\end{proof}

\begin{proposition}\label{b3:prop:differential-base-change}
设 \(B\) 为交换 \(A\) 代数，\(A\to A'\) 为交换环同态，令 \(B'=A'\otimes_AB\)。有自然同构
\[
\Omega_{B'/A'}\cong B'\otimes_B\Omega_{B/A}.
\]
不需要 \(A'\) 在 \(A\) 上平坦。
\end{proposition}
\begin{proof}
任取 \(B'\) 模 \(M\)，沿环同态 \(B\to B'\)、\(b\mapsto1\otimes b\) 将它看作 \(B\) 模。将 \(A'\) 导子 \(\Delta:B'\to M\) 与这个环同态复合，便给出 \(A\) 导子 \(D:B\to M\)。反过来，给定这样的 \(D\)，在纯张量上定义
\[
\Delta(a'\otimes b)=(a'\otimes1)\cdot D(b).
\]
由 \(D\) 的 \(A\) 线性及张量积中 \(A\) 的两个像相等，此公式保持平衡关系，因而对有限和良定义。它零化 \(A'\) 的像；对纯张量的乘积使用 \(D\) 的 Leibniz 法则，再由加法展开，就得到 \(\Delta\) 的 Leibniz 法则。复合与延拓互逆，因为任何 \(A'\) 导子都满足所列纯张量公式。因此两种导子自然一一对应。

又由普遍导子的泛性质，\(D\) 唯一写成 \(h\circ\mathrm{d}_{B/A}\)，其中 \(h:\Omega_{B/A}\to M\) 为 \(B\) 线性映射。这样的 \(h\) 又唯一延拓为 \(B'\) 线性映射
\[
B'\otimes_B\Omega_{B/A}\longrightarrow M,
\qquad b'\otimes\eta\longmapsto b'\cdot h(\eta).
\]
所以右侧张量积和 \(\Omega_{B'/A'}\) 对每个目标模都表示同一个导子函子，普遍对象的唯一性给出所需同构。这个同构将 \(1\otimes\mathrm{d}b\) 送到 \(\mathrm{d}(1\otimes b)\)。
\end{proof}

这两个公式允许先在多项式商上写出 Jacobian 展示，再局部化到指定点，或改变基域后重新考察方程。代数作为环的约化性和正则性可能在基变换后改变；光滑性则是相对于基环的性质，并在任意基变换下保持，具体论证见\secref{b3:sec:1405}。

% !TeX root = ../../Book3.tex
% 纲要标题：### 14.3 切空间与 Jacobian 矩阵
% 纲要位置：工作记录/计划纲要.md，第 2417 行。
% 纲要细目（写作范围）：
% * 极大理想模平方与余切空间
% * 仿射方程的切空间计算
% * 一般素点的剩余域微分方向；不可分闭点与基变换后的幂零结构

\section{切空间与 Jacobian 矩阵}
\label{b3:sec:1403}

一条曲线在一个点附近可能有几条分支，也可能含有重根；只看点集不能分辨这些现象。把点的坐标改成 \(a_i+v_i\varepsilon\)、其中 \(\varepsilon^2=0\)，方程留下的一次关系便限制了允许的方向。先从有理点说明这个计算，再讨论剩余域不是基域的情形。


\subsection{极大理想模平方与余切空间}
\label{b3:subsec:1403:01}

\begin{definition}\label{b3:def:zariski-tangent-space}
设 \(k\) 为域，\(B\) 为有限型交换 \(k\) 代数，\(a:B\to k\) 为一个 \(k\) 代数同态，\(\mathfrak m=\ker a\)。以 \(a\) 赋予 \(k\) 一个 \(B\) 模结构。向量空间
\[
T_a(B)=\operatorname{Der}_k(B,k)
\]
称为 \(B\) 在这个 \(k\) 有理点处的\term[zariski-qiekongjian]{Zariski 切空间}[Zariski tangent space]；\(\mathfrak m/\mathfrak m^2\) 是该点的余切空间。这里的 \(B\) 可以含有非零幂零元，不要求它是某个点集的约化坐标环。
\symidx{\(T_a(B)\)：有理点处的 Zariski 切空间}
\end{definition}

\begin{proposition}\label{b3:prop:tangent-cotangent-rational}
设 \(k\) 为域，\(B\) 为交换 \(k\) 代数，\(a:B\to k\) 为 \(k\) 代数同态，\(\mathfrak m=\ker a\)。则有自然同构
\[
\mathfrak m/\mathfrak m^2\xrightarrow{\sim}
k\otimes_B\Omega_{B/k},\qquad
\bar b\longmapsto1\otimes \mathrm{d}b,
\]
以及
\[
\operatorname{Der}_k(B,k)
\cong\operatorname{Hom}_k(\mathfrak m/\mathfrak m^2,k).
\]
若以局部环 \(B_{\mathfrak m}\) 代替 \(B\)，所得余切空间及切空间相同。
\end{proposition}
\begin{proof}
对 \(b,c\in\mathfrak m\)，\(1\otimes \mathrm{d}(bc)=0\)，故第一个映射良定义。每个 \(b\in B\) 都写成 \(a(b)+\bigl(b-a(b)\bigr)\)，其中后项属于 \(\mathfrak m\)，所以这个映射满射。反向定义
\[
D:B\longrightarrow\mathfrak m/\mathfrak m^2,
\qquad D(b)=b-a(b)\pmod{\mathfrak m^2}.
\]
将 \(b=a(b)+b_0\)、\(c=a(c)+c_0\) 相乘，模掉 \(\mathfrak m^2\) 后得到
\(D(bc)=a(b)\cdot D(c)+a(c)\cdot D(b)\)。故 \(D\) 是取值于剩余域模的 \(k\) 导子，由微分模泛性质诱导逆映射。第二个同构由同一泛性质及第一个同构得到。

若 \(s\notin\mathfrak m\)，它在 \(\mathfrak m/\mathfrak m^2\) 上的乘法就是非零标量 \(a(s)\)，已经可逆。因此局部化不改变该商；由局部化正合性，它又等于
\(\mathfrak mB_{\mathfrak m}/(\mathfrak mB_{\mathfrak m})^2\)。导子的分式延拓公式也给出相同的切空间。
\end{proof}

例如，在 \(k[x_1,\ldots,x_n]_{(x_1,\ldots,x_n)}\) 的原点，记极大理想为 \(\mathfrak m\)。消失于原点的函数模 \(\mathfrak m^2\) 后，只留下 \(\sum_i c_ix_i\) 这样的一次部分；\(\mathfrak m/\mathfrak m^2\) 保存的是函数的一阶线性信息。切向量 \(v=(v_i)\) 则是作用在这些一次部分上的线性泛函，将 \(\sum_i c_ix_i\) 送到 \(\sum_i c_iv_i\)。这就是切空间与余切空间互为对偶的具体含义。

上述导子还有直接的点的解释。给定 \(D\in T_a(B)\)，映射
\[
B\longrightarrow k[\varepsilon]/(\varepsilon^2),\qquad
b\longmapsto a(b)+D(b)\cdot\varepsilon
\]
是一个模 \(\varepsilon\) 后等于 \(a\) 的 \(k\) 代数同态；乘法相容性恰是导子的公式。反过来，每个这样的同态都唯一给出一个 \(D\)。因此，同一个切向量既可视为导子，也可视为余切空间上的线性泛函，或点向双数环中的一阶延伸。

\subsection{仿射方程的切空间计算}
\label{b3:subsec:1403:02}

\begin{theorem}\label{b3:thm:jacobian-tangent-space}
设 \(k\) 为域，\(B=k[X_1,\ldots,X_n]/(f_1,\ldots,f_s)\)，且 \(a=(a_1,\ldots,a_n)\in k^n\) 满足所有 \(f_i(a)=0\)。以 \(X_j\mapsto a_j\) 的求值同态 \(B\to k\) 定义该点的切空间 \(T_a(B)=\operatorname{Der}_k(B,k)\)。令
\[
J_f(a)=\left(\dfrac{\partial f_i}{\partial X_j}(a)\right)_{
1\le i\le s,\,1\le j\le n}.
\]
则在变量坐标给出的识别下，
\[
T_a(B)=\ker\bigl(J_f(a):k^n\to k^s\bigr),
\qquad \dim_kT_a(B)=n-\operatorname{rank}J_f(a).
\]
这个矩阵称为所给方程组的\term[jacobian-juzhen]{Jacobian 矩阵}[Jacobian matrix]，亦称雅可比矩阵\termidx[yakebi-juzhen]{雅可比矩阵}。\footnote{“雅可比矩阵”见菲赫金哥尔茨《微积分学教程》第一卷\cite[第六章“函数矩阵(雅可比矩阵)的乘法”]{Fichtenholz2005CalculusChineseI}。}
\end{theorem}
\begin{proof}
由双数值点的描述，一个方向 \(v=(v_j)\) 对应 \(X_j\mapsto a_j+v_j\varepsilon\)。它能从多项式环下降到 \(B\)，当且仅当每个 \(f_i\) 的像为零。因为 \(\varepsilon^2=0\)，二次以上的改变量全部消失，故
\[
f_i(a_1+v_1\varepsilon,\ldots,a_n+v_n\varepsilon)
=f_i(a)+\varepsilon\sum_j
\dfrac{\partial f_i}{\partial X_j}(a)v_j.
\]
由 \(f_i(a)=0\)，这一条件恰为 \(J_f(a)v=0\)，得到核的描述及维数公式。

\cref{b3:cor:differential-presentation}还将这一计算与环本身的微分模联系起来：展示 \(B^s\to B^n\to\Omega_{B/k}\to0\) 的关系矩阵为 \(J_f^{\mathsf t}\)。张量到 \(k\) 后仍右正合，再取 \(k\) 线性对偶，得到
\[
0\longrightarrow\operatorname{Hom}_k(k\otimes_B\Omega_{B/k},k)
\longrightarrow k^n\xrightarrow{J_f(a)}k^s.
\]
前一命题把左端自然识别为切空间。因此坐标与方程展示虽会改变矩阵，所得导子空间始终由 \(B\) 及这个有理点决定。
\end{proof}

例如设 \(k\) 为域，在 \(B=k[x,y]/(xy-1)\) 中，每个有理点 \((a,b)\) 都有 \(ab=1\)，Jacobian 行为 \((b,a)\)，秩为一，所以切空间由 \(bv_x+av_y=0\) 给出，是一维的。在尖点曲线 \(y^2=x^3\) 的原点，Jacobian 行 \((-3x^2,2y)\) 为零，故切空间是整个 \(k^2\)，任意特征都如此。曲线 \(xy=0\) 的原点也有二维切空间，虽然它有两条分支，而尖点只有一条。这说明切空间维数能发现某些异常，却不能单独数出分支。

\ProseParagraphBreak
尖点的参数化 \(t\mapsto(t^2,t^3)\) 把源直线原点的一阶点 \(t=v\varepsilon\) 送到 \(\bigl((v\varepsilon)^2,(v\varepsilon)^3\bigr)=(0,0)\)。因此它在原点诱导的切映射为零，源点的一阶变化在像中全部消失。

竖直方向则显示 Zariski 一阶切方向与切锥约化方向之间的差别。映射 \(x\mapsto0,y\mapsto\varepsilon\) 在 \(k[\varepsilon]/(\varepsilon^2)\) 中满足尖点方程，因此给出一个切向量。若提升到 \(k[\varepsilon]/(\varepsilon^3)\)，保持这一阶值就必须写成 \(x=a\varepsilon^2\)、\(y=\varepsilon+b\varepsilon^2\)。此时 \(x^3=0\)，而 \(y^2=\varepsilon^2\ne0\)，所以没有这样的提升。Zariski 切空间只检查一次关系；切锥 \(Y^2=0\) 则已保留最低二次关系，其约化支集只有水平直线。因而一阶允许的方向未必能延续成更高阶曲线。

\ProseParagraphBreak
还须保留方程中的幂零结构。\(k[t]/(t^2)\) 只有一个域值点，但该点切空间为一维；约化后的坐标环 \(k\) 则有零切空间。两者的点集相同，一阶延伸却不同。若已经把理想替换成根理想再计算，便无法恢复丢掉的方向。

\ProseParagraphBreak
Jacobian 计算还可以寻找参数族中切空间异常增大的位置。以 \(t\) 为参数，在特征零域 \(k\) 上考虑曲线族
\[
F_t(x,y)=y^2-x^3-t=0.
\]
固定参数后，切空间维数为二恰好要求 \(F_t=0\)、\(\partial F_t/\partial x=-3x^2=0\)、\(\partial F_t/\partial y=2y=0\)。在 \(k[t,x,y]\) 中，相应理想为
\[
(y^2-x^3-t,-3x^2,2y)=(y,x^2,t).
\]
右边的包含关系可直接由 \(t=y^2-x^3-F_t\) 核验，反向也由各生成元的表达式得到。这个 Jacobian 理想记录异常位置的代数结构，其商中还保留 \(x^2=0\) 的加厚信息；若只求底层异常点集，则取根理想得到 \((x,y,t)\)，只剩原点。要判断哪些参数会产生异常，还须消去 \(x,y\)，收缩到参数环得到
\[
(y,x^2,t)\cap k[t]=(t).
\]
这是因为 \(t\) 已属于该理想，而属于它的参数多项式必在原点取值为零。于是异常只在 \(t=0\) 出现；其余参数的每个几何点都有一维切空间。

\subsection{非完美基域上的切空间}
\label{b3:subsec:1403:03}

这里先比较一般素点处的两个一阶对象，差别并不限于不完美基域。设 \(k\) 为域，\(B\) 为交换 \(k\) 代数，\(\mathfrak p\in\operatorname{Spec}B\)，记 \(\kappa=\kappa(\mathfrak p)\)。剩余域 \(\kappa\) 未必等于 \(k\)。局部环自身的余切空间
\(\mathfrak pB_{\mathfrak p}/(\mathfrak pB_{\mathfrak p})^2\)
与相对微分的纤维 \(\kappa\otimes_B\Omega_{B/k}\) 出自不同构造。对 \(B_{\mathfrak p}\to\kappa\) 应用余法正合列及微分局部化，得到下列 \(\kappa\) 向量空间的右正合列：
\begin{equation}\label{b3:eq:cotangent-residue-sequence}
\begin{tikzcd}[column sep=small,row sep=0.8em]
\dfrac{\mathfrak pB_{\mathfrak p}}{(\mathfrak pB_{\mathfrak p})^2}
  \arrow[r]
&\kappa\otimes_B\Omega_{B/k}
  \arrow[r]
&\Omega_{\kappa/k}\arrow[r]&0\\
{\begin{gathered}\text{局部环自身}\\\text{的一阶结构}\end{gathered}}
&{\begin{gathered}\text{相对于基域 }k\\\text{的一阶结构}\end{gathered}}
&{\begin{gathered}\text{剩余域扩张}\\\text{的一阶结构}\end{gathered}}&
\end{tikzcd}
\end{equation}
在 \(k\) 有理点处，\cref{b3:prop:tangent-cotangent-rational}证明左箭头为同构。一般点处的左箭头不必单射，右端也可能非零；它记录了剩余域扩张本身贡献的微分。记左箭头的像为 \(W\)，则中间空间包含子空间 \(W\)，商空间为 \(\Omega_{\kappa/k}\)。向量空间的短正合列可以选基分裂，所以中间空间可非典范地写成 \(W\oplus\Omega_{\kappa/k}\)；但没有自然指定的补空间，而且 \(W\) 只是左端的一个商，不能直接把整个局部余切空间与 \(\Omega_{\kappa/k}\) 作直和。

即使 \(k\) 完美，在 \(B=k[t]\) 的一般点 \(\mathfrak p=(0)\) 处也有 \(\kappa=k(t)\)；由多项式微分与局部化公式，\(\Omega_{\kappa/k}=k(t)\,\mathrm{d}t\ne0\)。此时局部环是域，局部余切空间为零，相对微分纤维却仍有一维。非完美基域上的额外现象是：连闭点的有限剩余域扩张也可能贡献这样的非零微分。

\begin{example}\label{b3:exm:inseparable-tangent}
设 \(p\) 为素数、\(k=\mathbb F_p(s)\)，其中 \(s\) 为超越元。令
\[
L=k[u]/(u^p-s).
\]
由于 \(s\) 在 \(k\) 中不是 \(p\) 次幂，\(u^p-s\) 不可约，故 \(L\) 为域。其唯一局部环的极大理想为零，局部环余切空间也为零。但微分展示给出
\[
\Omega_{L/k}=L\,\mathrm{d}u,
\]
因为 \(s\) 是基域中的常数，且 \(k\) 的特征为 \(p\)，故 \(\mathrm{d}(u^p-s)=0\)。该点的剩余域就是 \(L\)，所以相对微分纤维 \(L\otimes_L\Omega_{L/k}=\Omega_{L/k}\) 非零，完全来自剩余域扩张的不可分性，并非来自局部极大理想。若扩域到 \(K=k(s^{1/p})\)，则
\[
K\otimes_kL\cong K[v]/(v^p),\qquad v=u-s^{1/p},
\]
其唯一 \(K\) 有理点由 \(v\mapsto0\) 给出，局部余切空间 \((v)/(v^2)\) 以 \(v\) 的类为基，故该点的切空间为一维。这里 \(v\) 的像是非零幂零元：原来的环 \(L\) 是域，改变基域后却出现了这个加厚点。
\end{example}

不可约性的一个直接检验如下。在代数闭包中，\(T^p-s=(T-\alpha)^p\)。若 \(\alpha\) 的最小多项式次数 \(r<p\)，它只能为 \((T-\alpha)^r\)，其 \(T^{r-1}\) 系数使 \(r\alpha\in k\)，从而 \(\alpha\in k\)，与 \(s\) 非 \(p\) 次幂矛盾。最后一事实可由有理函数在 \(s=0\) 处的零点阶数核验：\(p\) 次幂的阶数是 \(p\) 的倍数，\(s\) 的阶数则为一。

\ProseParagraphBreak
这一区别在不完美基域上尤其必要。局部环是域，因而在环论意义下最简单，却不能据此断言它相对于给定基域光滑。后面的正则局部环与光滑性判据将分别处理环自身的性质和扩张基域后的性质；本章的微分计算已经显示两者为何不能合并。

% !TeX root = ../../Book3.tex
% 纲要标题：### 14.4 域扩张、可分性与微分
% 纲要位置：工作记录/计划纲要.md，第 2399 行。
% 纲要细目（写作范围）：
% * 单个代数元的微分与有限域扩张的可分性判据
% * 指数一纯不可分扩张；子域 \(K L^p\) 及微分基的构造
% * 纯不可分扩张的次数与微分维数；基变换后的幂零结构

\section{域扩张、可分性与微分}
\label{b3:sec:1404}

代数元的关系 \(f(\alpha)=0\) 对导子给出 \(f'(\alpha)\mathrm{d}\alpha=0\)。若导数在域中非零，这个方向被迫消失；若关系是 \(p\) 次幂，求导却不能看见它。本节从单个代数元的计算出发，说明这种差别如何检测一般有限域扩张的可分性，再比较纯不可分扩张的次数与微分维数。

\subsection{有限代数域扩张的微分模}
\label{b3:subsec:1404:01}

先看一个代数元。设 \(K\) 为域，\(\alpha\) 为某个扩域中对 \(K\) 代数的元素，\(L=K(\alpha)\)，\(f\in K[T]\) 为 \(\alpha\) 的首一不可约多项式。由于 \(L=K[T]/(f)\)，\cref{b3:cor:differential-presentation}给出
\begin{equation}\label{b3:eq:simple-field-differentials}
\Omega_{L/K}\cong L\,\mathrm{d}\alpha/\bigl(f'(\alpha)\mathrm{d}\alpha\bigr).
\end{equation}
若 \(f\) 可分，\(f'(\alpha)\ne0\) 在域中为单位，故微分模为零。若 \(f\) 不可分，其不可约性迫使 \(f'=0\)，于是微分模为一维。一般有限扩张不一定由一个元素生成；不能把这个单生成元计算直接当作一般情形的证明。

\begin{lemma}\label{b3:lem:exponent-one-differentials}
设 \(K\subseteq L\) 为特征 \(p>0\) 的有限域扩张，令 \(F=K L^p\) 为由 \(K\) 及所有 \(p\) 次幂生成的子域。存在 \(\alpha_1,\ldots,\alpha_r\in L\)，使
\[
[L:F]=p^r,\qquad
\Omega_{L/K}\cong\bigoplus_{i=1}^r L\,\mathrm{d}\alpha_i.
\]
其中 \(r=0\) 恰好表示 \(L=K L^p\)。
\end{lemma}
\begin{proof}
每个 \(K\) 导子都零化 \(L^p\)，再由四则运算法则零化 \(F\)，所以 \(K\) 导子与 \(F\) 导子相同，微分模也相同。逐次选取尚未落入已生成子域的 \(\alpha_i\)，直至生成 \(L/F\)。记前一步生成的子域为 \(F_{i-1}=F(\alpha_1,\ldots,\alpha_{i-1})\)。因为 \(\alpha_i^p\in F\)，其在 \(F_{i-1}\) 上的最小多项式整除 \(T^p-\alpha_i^p\)。若次数为 \(j<p\)，则在代数闭包中必为
\[
(T-\alpha_i)^j=T^j-j\alpha_iT^{j-1}+\cdots,
\qquad 1\le j<p.
\]
系数 \(-j\alpha_i\) 属于 \(F_{i-1}\)，而 \(j\) 在特征 \(p\) 的域中非零，故 \(\alpha_i\in F_{i-1}\)，与选取方式矛盾。因此最小多项式恰为 \(T^p-\alpha_i^p\)，每步次数为 \(p\)，而有限次数保证过程终止。

于是 \(\prod_i\alpha_i^{e_i}\)、\(0\le e_i<p\)，组成 \(F\) 基。自然满射
\[
F[X_1,\ldots,X_r]/(X_1^p-\alpha_1^p,\ldots,X_r^p-\alpha_r^p)
\longrightarrow L
\]
两侧维数都是 \(p^r\)，故为同构。关系的微分全部为零，\cref{b3:cor:differential-presentation}遂给出所列自由基。
\end{proof}

\begin{theorem}\label{b3:thm:finite-field-differentials-separable}
设 \(L/K\) 为有限域扩张。则
\[
\Omega_{L/K}=0\quad\Longleftrightarrow\quad L/K\text{ 可分}.
\]
若 \(L/K\) 为特征 \(p>0\) 的非平凡有限纯不可分扩张，则 \(\Omega_{L/K}\ne0\)。
\end{theorem}
\begin{proof}
若扩张可分，对任意 \(\alpha\in L\)，其最小多项式 \(f\) 满足 \(f'(\alpha)\ne0\)。微分等式 \(f'(\alpha)\mathrm{d}\alpha=0\) 因而给出 \(\mathrm{d}\alpha=0\)。这些元素生成微分模，所以该模为零。

先证非平凡纯不可分情形。设 \(L/K\) 的指数不超过 \(e\)，即 \(L^{p^e}\subseteq K\)；有限生成性保证存在同一个 \(e\)。若 \(L=K L^p\)，反复取 \(p\) 次幂并代回便得
\(L=K L^{p^j}\) 对所有 \(j\ge1\) 成立，特别地 \(L=K\)，矛盾。因此 \(K L^p\subsetneq L\)，前一引理给出至少一个非零微分。

最后设 \(L/K\) 不可分。令 \(E\) 为 \(L\) 中对 \(K\) 可分的元素组成的子域；可分代数扩张对复合及子扩张封闭，保证这些元素对四则运算封闭。\(E/K\) 有限可分。对任意 \(\alpha\in L\)，其不可约多项式在特征 \(p\) 下唯一写成 \(g(T^{p^a})\)，其中 \(g'\ne0\)。多项式 \(g\) 仍不可约，否则将其分解代入 \(T^{p^a}\) 就得到原多项式的分解；故 \(g\) 可分，\(\alpha^{p^a}\) 对 \(K\) 可分并属于 \(E\)。于是 \(L/E\) 纯不可分；由于原扩张不可分，它非平凡。复合微分正合列给出
\[
L\otimes_E\Omega_{E/K}\longrightarrow\Omega_{L/K}
\longrightarrow\Omega_{L/E}\longrightarrow0.
\]
左端由已证的可分情形为零，右端由刚证的纯不可分情形非零，所以中间项非零。特征零没有不可分代数扩张，不需此分情形。
\end{proof}

\subsection{纯不可分扩张的微分计算}
\label{b3:subsec:1404:05}

微分维数记录独立的一阶方向，却未必记录纯不可分扩张的全部幂次深度。设 \(p\) 为素数，\(s\) 为超越元，\(k=\mathbb F_p(s)\)、\(L=k(\alpha)\)，其中 \(\alpha^{p^e}=s\)、\(e\ge1\) 为整数。则
\[
\Omega_{L/k}=L\,\mathrm{d}\alpha.
\]
最小多项式 \(T^{p^e}-s\) 的不可约性可用 \(s\) 在 \(\mathbb F_p[s]\) 中的 Eisenstein 判据检验。扩张次数为 \(p^e\)，微分维数却一直为一，因此微分维数不是不可分次数的对数。它测量的是独立的一阶方向，未记录每个方向的全部幂次深度。

\ProseParagraphBreak

多方向的例子可由两个代数独立的元 \(s,t\) 构造。设 \(p\) 为素数，取
\[
k=\mathbb F_p(s,t),\qquad
L=k(\alpha,\beta),\qquad \alpha^p=s,\quad\beta^p=t.
\]
有 \([L:k]=p^2\) 及
\(\Omega_{L/k}=L\,\mathrm{d}\alpha\oplus L\,\mathrm{d}\beta\)。这里可把 \(L\) 识别为有理函数域 \(\mathbb F_p(\alpha,\beta)\)：不同模 \(p\) 的指数类在多项式恒等式中不会相消，所以 \(\alpha^i\beta^j\)、\(0\le i,j<p\)，对 \(k=\mathbb F_p(\alpha^p,\beta^p)\) 线性无关，并显然生成。对应的两个关系微分都为零，给出所列基。

\ProseParagraphBreak
在这一例子中，原来的 \(L\) 是域，扩到 \(L\) 本身后却有
\[
L\otimes_kL\cong L[\varepsilon,\eta]/(\varepsilon^p,\eta^p).
\]
把右边变量分别送到第二因子的 \(\alpha,\beta\) 减去第一因子的相应元素，便得到满射；两侧 \(L\) 维数都是 \(p^2\)，所以为同构。与\cref{b3:exm:inseparable-tangent}一样，这个计算把不可分性表现为基变换后出现的一阶方向和幂零结构。

% !TeX root = ../../Book3.tex
% 纲要标题：### 14.5 形式光滑、非分歧与平展
% 纲要位置：工作记录/计划纲要.md，第 2405 行。
% 纲要细目（写作范围）：
% * 平方零提升的存在性、至多唯一性与存在唯一性
% * 六种无穷小性质的条件对照；非分歧采用有限型，光滑和平展采用有限展示
% * 微分模消失与形式非分歧的等价
% * 标准平展的 Newton 修正与 Jacobian 子式可逆时的光滑提升
% * 形式及普通性质的任意基变换；有限域扩张平展与可分性的等价
% * 一般有限展示形式平展的等价刻画及光滑蕴含平坦在第23章证明
% ---

\section{形式光滑、非分歧与平展}
\label{b3:sec:1405}

导子比较的是两个同态之间的一阶差别；给定商环中的一个同态，能否把它提升到扩张环中，却是另一个问题。有限域扩张的微分计算还显示，环本身是域也不保证它相对于基域具有良好的提升性质。本节将存在性与唯一性分别表述，再用导数可逆的方程构造提升，并解释有限性和平坦性在光滑、非分歧及平展条件中的作用。

\subsection{平方零提升与基本概念}
\label{b3:subsec:1405:01}
\label{b3:subsec:1404:02}

平方零扩张已经把导子解释为同态的一阶变化。现在还要问：商环中的一个同态，是否来自扩张环中的同态？若 \(B=A[X_1,\ldots,X_n]/(f_1,\ldots,f_s)\)，任取变量像的代表后，各 \(f_j\) 的值都落在被商去的理想中。导子控制不同修正之间的差，但关系中的误差未必能够消掉。因此，提升的存在性和唯一性是两个需要分别考虑的问题。

\begin{definition}\label{b3:def:formal-infinitesimal-properties}
设 \(A\to B\) 为交换环同态。对任意交换 \(A\) 代数 \(C\)、满足 \(J^2=0\) 的理想 \(J\subseteq C\)，以及 \(A\) 代数同态 \(u:B\to C/J\)，考虑使下式复合等于 \(u\) 的 \(A\) 代数同态 \(\widetilde u\)：
\[
B\xrightarrow{\widetilde u}C\longrightarrow C/J.
\]
若每组这样的数据都有提升，则称 \(A\to B\) \term[xingshi-guanghua]{形式光滑}[formally smooth]；若每组数据至多有一个提升，则称其\term[xingshi-feifenqi]{形式非分歧}[formally unramified]；若每组数据恰有一个提升，则称其\term[xingshi-pingzhan]{形式平展}[formally étale]。
\end{definition}

“至多一个”允许没有提升，“至少一个”允许有多个，必须分开判断。只用平方零理想已足以检验任意幂零理想：若 \(J^N=0\)，可沿 \(C/J\)、\(C/J^2\)、\ldots、\(C/J^N=C\) 逐次提升，因为 \(J^i/J^{i+1}\) 的平方为零；存在性、唯一性分别逐次传递。

\ProseParagraphBreak
例如设 \(k\) 为域，取
\[
B=k[z]/(z^2),\qquad C=k[\varepsilon]/(\varepsilon^3),
\qquad J=(\varepsilon^2).
\]
在 \(C/J\) 中，\(z\mapsto\varepsilon\) 保持关系，因而给出 \(k\) 代数同态。若它能够提升，则 \(z\) 的像必为 \(\varepsilon+a\varepsilon^2\)，其中 \(a\in k\)。但在 \(C\) 中
\[
(\varepsilon+a\varepsilon^2)^2=\varepsilon^2\ne0,
\]
无论怎样选 \(a\)，误差都不能消去。因此 \(B\) 在 \(k\) 上不形式光滑。若去掉关系，改用 \(k[z]\)，同样的近似同态却有全部 \(z\mapsto\varepsilon+a\varepsilon^2\) 这些提升；自由变量使提升存在，也留下不唯一的变化方向。

\begin{definition}\label{b3:def:finite-presentation-infinitesimal-properties}
设 \(A\to B\) 为交换环同态。若存在有限个 \(b_1,\ldots,b_n\in B\)，使 \(B=A[b_1,\ldots,b_n]\)，则称 \(B\) 为有限型 \(A\) 代数。若 \(B\) 还能写成
\(A[X_1,\ldots,X_n]/(f_1,\ldots,f_s)\)，其中 \(n,s\) 为非负整数，则称 \(B\) 为有限展示 \(A\) 代数（也称有限呈示 \(A\) 代数）。对同态 \(A\to B\)，分别作如下定义：若 \(B\) 有限展示且该同态形式光滑，则称其\term[guanghua]{光滑}[smooth]；若 \(B\) 有限型且 \(\Omega_{B/A}=0\)，则称其\term[feifenqi]{非分歧}[unramified]；若 \(B\) 有限展示、作为 \(A\) 模平坦且 \(\Omega_{B/A}=0\)，则称其\term[pingzhan]{平展}[étale]。
\end{definition}

非分歧只要求有限型；光滑和平展的定义则要求有限展示。由下文的\cref{b3:thm:formally-unramified-differentials}，非分歧也等价于有限型且形式非分歧。\cref{b3:thm:etale-equivalent-characterizations}将证明有限展示且形式平展等价于平展，\cref{b3:thm:smooth-flat}将证明光滑蕴含平坦。六种性质的条件列于\cref{b3:tab:infinitesimal-properties}。

\begin{center}
\begin{minipage}{\textwidth}
\makeatletter
\def\@captype{table}
\makeatother
\centering
\caption[无穷小提升性质的条件]{无穷小提升性质的条件。固定交换环同态 \(A\to B\)；提升栏对任意交换 \(A\) 代数 \(C\)、平方零理想 \(J\subseteq C\) 及 \(A\) 代数同态 \(B\to C/J\) 而言。平展行的提升刻画及光滑的平坦性由第23章的相应定理给出。}
\label{b3:tab:infinitesimal-properties}
\begin{tabularx}{\textwidth}{@{}>{\raggedright\arraybackslash}p{0.17\textwidth}>{\raggedright\arraybackslash}p{0.18\textwidth}>{\raggedright\arraybackslash}p{0.27\textwidth}>{\raggedright\arraybackslash}X@{}}
\toprule
性质 & 平方零提升 & 有限性条件 & 平坦性的地位\\
\midrule
形式光滑 & 至少一个 & 不要求有限性 & 定义不要求\\
形式非分歧 & 至多一个 & 不要求有限性 & 定义不要求\\
形式平展 & 恰好一个 & 不要求有限性 & 定义不要求\\
光滑 & 至少一个 & 有限展示 & 可由定理推出\\
非分歧 & 至多一个 & 有限型 & 不要求\\
平展 & 恰好一个 & 有限展示 & 定义要求\\
\bottomrule
\end{tabularx}
\end{minipage}
\end{center}

\ProseParagraphBreak
例如多项式代数 \(A[X_1,\ldots,X_n]\) 光滑：把每个 \(u(X_i)\) 任意提升到 \(C\)，即可唯一按这些已选的像延拓为代数同态。若 \(n>0\)，选择本身通常不唯一，所以不能据此称形式非分歧。局部化 \(A_f\) 则形式平展：\(f\) 在 \(C/J\) 中的像为单位时，它在 \(C\) 中也是单位；若 \(ab=1+j\)、\(j^2=0\)，则 \(b(1-j)\) 就是 \(a\) 的逆。局部化的泛性质因而给出唯一提升。再用 \(A_f=A[T]/(fT-1)\)、局部化平坦性及 \(\Omega_{A_f/A}=0\)，也可直接核验其平展性。

\subsection{微分模消失与提升唯一性}
\label{b3:subsec:1405:02}
\label{b3:subsec:1404:03}

\begin{theorem}\label{b3:thm:formally-unramified-differentials}
设 \(A\to B\) 为任意交换环同态，不加有限性条件。该映射形式非分歧，当且仅当 \(\Omega_{B/A}=0\)。
\end{theorem}
\begin{proof}
设 \(J^2=0\)，且 \(v,w:B\to C\) 为同一个 \(u:B\to C/J\) 的两个提升。差 \(D=v-w\) 取值于 \(J\)。由于 \(J^2=0\)，\(u\) 使 \(J\) 成为一个 \(B\) 模：取 \(u(b)\) 的任意提升乘 \(J\)，结果不依赖选择。展开
\[
v(b)v(c)-w(b)w(c)
=w(b)D(c)+w(c)D(b)+D(b)D(c)
\]
可见最后一项为零，故 \(D\) 是 \(A\) 导子。如果 \(\Omega_{B/A}=0\)，微分模泛性质使 \(D=0\)，两个提升相同。

反向在平方零扩张 \(C=B\oplus\Omega_{B/A}\) 中比较
\(b\mapsto(b,0)\) 与 \(b\mapsto(b,\mathrm{d}b)\)。\cref{b3:prop:derivation-square-zero}说明两者都是恒等同态 \(B\to B\) 的提升。形式非分歧使它们相同，故所有 \(\mathrm{d}b=0\)。这些元素生成 \(\Omega_{B/A}\)，所以该模为零。
\end{proof}

这项等价也解释了微分模的作用：它的消失控制两个已存在的提升之间的差；提升的存在性与平坦性仍须分别检验。

\begin{example}\label{b3:exm:unramified-not-etale}
设 \(k\) 为域，取商映射 \(A=k[t]\to B=k=A/(t)\)。任意 \(A\) 导子 \(B\to M\) 都为零，因为 \(B\) 的每个元素都来自基环，因此 \(\Omega_{B/A}=0\)，且该映射有限展示、形式非分歧。但它不平坦：单射 \(A\xrightarrow{t}A\) 张量到 \(B\) 后变成非零模 \(B\) 上的零映射。

\ProseParagraphBreak
它也不是形式光滑。把 \(C=k[t]/(t^2)\) 看成 \(A\) 代数，令 \(J=(t)/(t^2)\)。恒等映射 \(B=k\to C/J=k\) 不可能提升为 \(A\) 代数映射 \(B\to C\)，因为 \(t\) 在 \(B\) 中为零，在 \(C\) 中却非零。故它是非分歧而非平展的具体例子。
\end{example}

\subsection{标准平展与 Jacobian 光滑代数}
\label{b3:subsec:1405:03}
\label{b3:subsec:1404:04}

微分判据说明非分歧控制提升的唯一性。要得到提升，还需消去关系中的误差；标准平展代数中，导数可逆使这一步成为一次线性修正。

\begin{proposition}\label{b3:prop:standard-etale-lift}
设 \(A\) 为交换环，\(f\in A[T]\) 为正次数首一多项式，令
\(B=\bigl(A[T]/(f)\bigr)_g\)，其中 \(g\) 是商环中的一个元素，\(t\) 为 \(T\) 的像。若 \(f'(t)\) 在 \(B\) 中可逆，则 \(A\to B\) 平展，并且形式平展。
\end{proposition}
\begin{proof}
商环由首一除法具有有限 \(A\) 基 \(1,t,\ldots,t^{\deg f-1}\)，因而平坦；再局部化仍对 \(A\) 平坦，由\cref{b3:prop:flat-base-change-transitivity}成立。添加一个生成元及关系 \(gU-1\) 表明 \(B\) 有限展示，而
\(\Omega_{B/A}=B\,\mathrm{d}t/\bigl(f'(t)\mathrm{d}t\bigr)=0\)。所以映射平展。

直接检查提升。设 \(J^2=0\)，给定 \(u:B\to C/J\)。取 \(x\in C\) 提升 \(u(t)\)，则 \(f(x)\in J\)，且 \(f'(x)\) 为单位，因为它模 \(J\) 后为单位。作一次 Newton 修正\footnote{牛顿（Isaac Newton，1643-1727），英国数学家及物理学家，创立微积分并提出经典力学定律。}，令
\[
x'=x-f'(x)^{-1}f(x).
\]
平方零条件使 Taylor 展开\footnote{泰勒（Brook Taylor，1685-1731），英国数学家，以函数的级数展开研究微积分。}恰好为
\(f(x')=f(x)+f'(x)(x'-x)=0\)。选取代表 \(g\) 的多项式 \(G\in A[T]\)。因为 \(G(x')\) 模 \(J\) 为单位，它本身为单位；而 \(f(x')=0\) 保证这个值不依赖代表的选择。于是 \(T\mapsto x'\) 唯一延拓为 \(B\to C\)。若另一根 \(x''\) 给出同一模 \(J\) 的像，则
\(0=f(x'')-f(x')=f'(x')(x''-x')\)，故 \(x''=x'\)。由生成元及局部化泛性质，整个提升唯一。
\end{proof}

这种一元计算通常称为标准平展代数的基本情形。条件必须在实际局部化的环内检查。例如
\(A=k[s]\)、\(B=A[t]/(t^2-s)\) 总是秩二自由，但微分模为 \(\bigl(B/(2t)\bigr)\,\mathrm{d}t\)。若 \(\operatorname{char}k\ne2\)，倒置 \(t\) 后导数为单位，得到平展映射；在 \(t=0\) 处导数不能相除。若 \(\operatorname{char}k=2\)，导数处处为零，非零的局部化也不满足这个非分歧条件。

光滑的类似计算允许自由方向：只要选出足够多的变量消去关系误差，其余变量的像仍可自由选择。

\begin{proposition}\label{b3:prop:jacobian-smooth-lifting}
设 \(A\) 为交换环，\(1\le r\le n\) 为整数，\(f_1,\ldots,f_r\in A[X_1,\ldots,X_n]\)，\(g\) 为商环 \(A[X_1,\ldots,X_n]/(f_1,\ldots,f_r)\) 的元素，记
\[
B=\bigl(A[X_1,\ldots,X_n]/(f_1,\ldots,f_r)\bigr)_g.
\]
若 Jacobian 矩阵 \(\bigl(\partial f_i/\partial X_j\bigr)\) 的某个 \(r\times r\) 子式在 \(B\) 中可逆，则 \(A\to B\) 光滑。
\end{proposition}
\begin{proof}
重新排列变量，可设前 \(r\) 列的子式可逆。给定交换 \(A\) 代数 \(C\)、平方零理想 \(J\subseteq C\) 及 \(A\) 代数同态 \(u:B\to C/J\)，先任取所有变量像的提升 \(x_j\in C\)。误差向量 \(\bigl(f_i(x)\bigr)_{i=1}^r\) 的各分量落入 \(J\)。所选子式在 \(C/J\) 中为单位，因而在 \(C\) 中仍为单位。保持后 \(n-r\) 个变量不变，用可逆的 \(r\times r\) 子矩阵解线性方程
\[
\left(\dfrac{\partial f_i}{\partial X_j}(x)\right)_{1\le i,j\le r}
\begin{pmatrix}h_1\\ \vdots\\h_r\end{pmatrix}
=-\begin{pmatrix}f_1(x)\\ \vdots\\f_r(x)\end{pmatrix}.
\]
所得 \(h_j\) 都属于 \(J\)。平方零条件给出准确的一阶等式
\[
\begin{aligned}
&f_i(x_1+h_1,\ldots,x_r+h_r,x_{r+1},\ldots,x_n)\\
&\qquad=f_i(x)+\sum_{j=1}^r
  \dfrac{\partial f_i}{\partial X_j}(x)h_j=0.
\end{aligned}
\]
因为 \(J^2=0\)，展开中含有至少两个 \(h_j\) 的项全部为零；最后一个等式正是所解的线性方程。\(g\) 的值模 \(J\) 后仍为单位，故修正后的值也可逆。于是修正后的变量像满足全部关系，并允许 \(g\) 的逆元，从而给出提升 \(B\to C\)。这证明 \(A\to B\) 形式光滑；添加一个生成元及关系 \(gU-1\) 可知 \(B\) 有限展示，故该同态光滑。
\end{proof}

\subsection{基变换与有限域扩张}
\label{b3:subsec:1405:04}

平方零提升的条件也能直接随基环扩张传递。设 \(A\to B\) 为交换环同态，\(A'\) 为交换 \(A\) 代数，记 \(B'=A'\otimes_AB\)。对任意 \(A'\) 代数 \(C\) 及平方零理想 \(J\)，映射 \(B'\to C/J\) 等价于一个 \(A\) 代数映射 \(B\to C/J\)。它的每个提升 \(B\to C\)，与固定的结构映射 \(A'\to C\) 在 \(A\) 上相同，故由张量积的泛性质唯一合成 \(B'\to C\)；反过来，限制到 \(B\) 还原该提升。因此两边的提升集合一一对应，存在性和至多唯一性分别保持。有限型、有限展示及平坦性也保持基变换，所以光滑、非分歧和平展都保持任意基变换。

\begin{corollary}\label{b3:cor:finite-field-extension-etale}
设 \(L/K\) 为有限域扩张。则 \(K\to L\) 非分歧，当且仅当它平展，也当且仅当 \(L/K\) 可分；此时 \(K\to L\) 还形式平展。
\end{corollary}
\begin{proof}
向量空间自由，故 \(L\) 作为 \(K\) 模平坦。有限扩张也是有限展示 \(K\) 代数：选有限组代数生成元，得到多项式环到 \(L\) 的满射，关系理想由 Hilbert 基定理有限生成。于是非分歧和平展都等价于 \(\Omega_{L/K}=0\)，再用\cref{b3:thm:finite-field-differentials-separable}得到可分性判据。若扩张可分，本原元定理给出 \(L=K(\alpha)\)；其最小多项式的导数在 \(L\) 中可逆，\cref{b3:prop:standard-etale-lift}便给出形式平展性。
\end{proof}

可分扩张的提升性质也可对照松村英之的《Commutative Ring Theory》\cite[\S25, Theorem 25.3, pp. 195--196]{Matsumura1989CommutativeRingTheory}。纯不可分扩张则可能使提升连存在性都失去。设 \(p\) 为素数，\(k=\mathbb F_p(s)\)，其中 \(s\) 为超越元；取 \(L=k[u]/(u^p-s)\)。置 \(f=T^p-s\)、\(C=k[T]/(f^2)\)、\(J=(f)/(f^2)\)。有 \(J^2=0\)、\(C/J=L\)。若恒等同态 \(L\to C/J\) 有提升，\(u\) 的像必为 \(T+h\)、\(h\in J\)。但
\[
(T+h)^p-s=T^p-s+h^p=f\ne0\quad\text{在 }C\text{ 中},
\]
因为 \(p\ge2\) 使 \(h^p=0\)，而 \(f\notin(f^2)\)。这与 \(u^p=s\) 矛盾，所以 \(k\to L\) 不形式光滑。


\begin{chaptersummary}
导子对乘积的一阶变化满足 Leibniz 法则\footnote{莱布尼茨（Gottfried Wilhelm Leibniz，1646-1716），德国数学家及哲学家，创立微积分，并研究符号演算。}，微分模把这些变化统一为一个泛性质。多项式环的微分自由，添加关系后则以关系的偏导数形成展示矩阵。余法正合列与复合微分正合列都只保证右正合，局部化和基变换公式也不替代其他性质所需的条件。

\ProseParagraphBreak
在有理点处，微分模的剩余纤维是极大理想模平方，切空间是它的对偶，也就是 Jacobian 矩阵的核。一般剩余域还会贡献自身的相对微分，不可分闭点表明两种余切对象可能不同。对有限域扩张，微分为零恰好等价于可分；纯不可分扩张中的非零微分可用指数为一的子扩张明确计算。

\ProseParagraphBreak
微分模 \(\Omega_{B/A}\) 消失，恰好表示每个平方零提升问题至多有一个提升；形式光滑则要求每个问题至少有一个提升。非分歧在前者之外要求有限型，光滑在后者之外要求有限展示。平展把有限展示、平坦性与微分消失合在一起；有限展示且形式平展与这些条件的等价将在\chapref{b3:ch:23}证明。标准平展代数的一次 Newton 修正、Jacobian 子式可逆时的线性修正，以及纯不可分扩张的提升障碍，分别展示了这些条件怎样作用于实际方程。
\end{chaptersummary}

\BookExercises

\begin{exercise}\label{b3:ex:14-derivation-constants}
设 \(k\) 为域、\(D:k[x]\to k[x]\) 为形式求导。求 \(\ker D\)，分别讨论特征零与特征 \(p>0\)。再求 \(k[x,y]\) 上同时被两个偏导子零化的多项式。
\end{exercise}
\begin{solution}
写 \(f=\sum_na_nx^n\)，则 \(Df=\sum_{n\ge1}na_nx^{n-1}\)。特征零时核为 \(k\)，特征 \(p\) 时只有指数被 \(p\) 整除的项可以保留，故核为 \(k[x^p]\)。两个变量的情形逐项比较系数，得到特征零时为 \(k\)，特征 \(p\) 时为 \(k[x^p,y^p]\)。这里系数可来自不完美域，不能把 \(k[x^p]\) 无条件写成 \(k[x]\) 中所有 \(p\) 次幂的集合。
\end{solution}

\begin{exercise}\label{b3:ex:14-euler-derivation}
设 \(k\) 为域，\(n\ge1\) 为整数。在 \(P=k[x_1,\ldots,x_n]\) 上令 \(E=\sum_i x_i\partial/\partial x_i\)。证明每个 \(d\) 次齐次多项式满足 \(E(f)=df\)。说明正特征时 \(E(f)=0\) 为什么不能推出 \(f\) 为常数。
\end{exercise}
\begin{solution}
对 \(x^\alpha\)，有 \(E(x^\alpha)=(\sum_i\alpha_i)x^\alpha\)。同次齐次项都有 \(\sum_i\alpha_i=d\)，线性相加得到公式。在特征 \(p\) 时，任意正次数为 \(p\) 倍数的齐次多项式都被零化，例如非常数的 \(x_1^p\)。当 \(n\ge2\) 时，\(x_1^{p-1}x_2\) 也被零化，且它并不是 \(p\) 次幂。
\end{solution}

\begin{exercise}\label{b3:ex:14-lifts-affine-space}
设 \(k\) 为域，\(B=k[x,y]/(xy)\)、\(C=k[\varepsilon]/(\varepsilon^3)\)、\(J=(\varepsilon^2)\subset C\)。对 \(a,b\in k\)，令 \(u:B\to C/J\) 将 \(x,y\) 分别送到 \(a\varepsilon,b\varepsilon\)。判定哪些 \(a,b\) 使 \(u\) 有提升，并求出全部提升。计算以 \(u\) 赋予 \(J\) 模结构后的 \(\operatorname{Der}_k(B,J)\)，说明为什么这个导子模的非零性不能保证提升存在。
\end{exercise}
\begin{solution}
两个生成元像的任意提升都写成
\[
x\longmapsto a\varepsilon+c\varepsilon^2,\qquad
y\longmapsto b\varepsilon+d\varepsilon^2,
\quad c,d\in k.
\]
它们的乘积为 \(ab\varepsilon^2\)，不受 \(c,d\) 影响。因此 \(ab=0\) 时每一对 \(c,d\) 都给出提升；\(ab\ne0\) 时没有提升。

\(x,y\) 对 \(J=k\varepsilon^2\) 的作用都为零，所以关系的微分 \(yD(x)+xD(y)=0\) 自动成立。两个导子值可在 \(J\) 中独立选择，得到 \(\operatorname{Der}_k(B,J)\cong J\oplus J\cong k^2\)，无论 \(a,b\) 如何都是如此。有初始提升时，这两个参数描述全部修正；没有初始提升时，它们却不能消去 \(ab\varepsilon^2\) 这个关系误差。
\end{solution}

\begin{exercise}\label{b3:ex:14-truncated-differentials}
设 \(k\) 为域、\(n\ge2\)、\(B=k[t]/(t^n)\)。计算 \(\Omega_{B/k}\) 的 \(k\) 维数，并计算余法映射 \((t^n)/(t^{2n})\to B\,\mathrm{d}t\) 的核。
\end{exercise}
\begin{solution}
微分展示为 \(\Omega_{B/k}=B\,\mathrm{d}t/(nt^{n-1}\cdot\mathrm{d}t)\)。若 \(\operatorname{char}k=0\)，或 \(\operatorname{char}k=p>0\) 且 \(p\nmid n\)，则 \(n\) 在 \(k\) 中非零，维数为 \(n-1\)；若 \(\operatorname{char}k=p>0\) 且 \(p\mid n\)，则维数为 \(n\)。余法模以 \(t^n\) 的类为一个 \(B\) 自由基，映射等同于乘 \(nt^{n-1}\)。在前一种情形中，其核由 \(t\) 生成；在后一种情形中，整个余法模都是核。两种情形中左箭头都不单射。
\end{solution}

\begin{exercise}\label{b3:ex:14-laurent-differentials}
设 \(k\) 为域，\(B=k[x^{\pm1},y^{\pm1}]\)。证明 \(\mathrm{d}x/x,\mathrm{d}y/y\) 为 \(\Omega_{B/k}\) 的自由基，计算任意整数 \(a,b\) 对应的 \(\mathrm{d}(x^ay^b)\)。
\end{exercise}
\begin{solution}
多项式微分局部化后，\(\mathrm{d}x,\mathrm{d}y\) 仍为自由基；分别乘单位 \(x^{-1},y^{-1}\) 得到所述新基。由 \(\mathrm{d}(x^{-1})=-x^{-2}\cdot\mathrm{d}x\) 及整数幂的归纳，
\[
\mathrm{d}(x^ay^b)=x^ay^b\left(a\dfrac{\mathrm{d}x}{x}+b\dfrac{\mathrm{d}y}{y}\right).
\]
公式中的整数解释为其在 \(k\) 中的像，因此正特征时仍有效。
\end{solution}

\begin{exercise}\label{b3:ex:14-tensor-product-differentials}
设 \(B,C\) 为交换 \(A\) 代数，\(D=B\otimes_A C\)。证明
\[
\Omega_{D/A}\cong
(D\otimes_B\Omega_{B/A})\oplus(D\otimes_C\Omega_{C/A}).
\]
\end{exercise}
\begin{solution}
取任意 \(D\) 模 \(M\)。一个 \(A\) 导子 \(\delta:D\to M\) 限制到 \(B\otimes1\)、\(1\otimes C\)，给出两个导子 \(\delta_B,\delta_C\)。它被这两者决定，因为
\[
\delta(b\otimes c)=(1\otimes c)\delta_B(b)+(b\otimes1)\delta_C(c).
\]
反过来，此公式保持 \(A\) 平衡关系，且乘法展开验证 Leibniz 法则，所以每一对导子都唯一延拓。右侧直和代表这两个导子的自由选择，因而具有 \(\Omega_{D/A}\) 的泛性质，得到同构。
\end{solution}

\begin{exercise}\label{b3:ex:14-power-map-differentials}
设 \(k\) 为域，\(m\ge2\)，环同态 \(k[u]\to k[t]\) 由 \(u\mapsto t^m\) 给定。计算 \(\Omega_{k[t]/k[u]}\)，并判断在何处局部化后微分模为零。
\end{exercise}
\begin{solution}
把 \(k[t]\) 写成 \(k[u][T]/(T^m-u)\)，得到
\(\Omega_{k[t]/k[u]}=k[t]\,\mathrm{d}t/(m t^{m-1}\cdot\mathrm{d}t)\)。若 \(\operatorname{char}k=0\)，或 \(\operatorname{char}k=p>0\) 且 \(p\nmid m\)，则微分模在素理想 \(\mathfrak p\) 处为零当且仅当 \(t\notin\mathfrak p\)：此时系数为单位；若 \(t\in\mathfrak p\)，系数落入局部极大理想，余模非零。若 \(\operatorname{char}k=p>0\) 且 \(p\mid m\)，则微分模自由秩一，在每个素理想处均非零。
\end{solution}

\begin{exercise}\label{b3:ex:14-cusp-torsion}
设 \(k\) 为域且 \(\operatorname{char}k\ne2,3\)，\(B=k[x,y]/(y^2-x^3)\)。正文已证明 \(\omega=2x\cdot\mathrm{d}y-3y\cdot\mathrm{d}x\) 是 \(\Omega_{B/k}\) 中的非零挠元。求它的完整零化理想 \(\operatorname{Ann}_B(\omega)\)，并求循环子模 \(B\omega\) 的 \(k\) 维数。
\end{exercise}
\begin{solution}
由微分展示，\(h\omega=0\) 当且仅当存在 \(q\in B\) 使
\[
(-3hy,2hx)=q(-3x^2,2y).
\]
将 \(B\) 识别为 \(k[t^2,t^3]\)，在分式域中比较第二坐标得 \(q=h/t\)；此值也满足第一坐标。因此 \(h\) 属于零化理想，恰好是 \(h/t\in B\)。\(B\) 中的多项式不含一次项，除以 \(t\) 后仍属于 \(B\)，当且仅当 \(h\) 的常数项与二次项都为零。这些多项式组成理想 \((t^3,t^4)=(y,x^2)\)，故
\[
\operatorname{Ann}_B(\omega)=(y,x^2),\qquad
B\omega\cong B/(y,x^2)\cong k[x]/(x^2).
\]
所以 \(B\omega\) 的 \(k\) 维数为二，基可取 \(\omega,x\omega\)。
\end{solution}

\begin{exercise}\label{b3:ex:14-parabola-characteristic-two}
设 \(k\) 为特征二的域，\(B=k[x,y]/(y-x^2)\)。计算每个 \(k\) 有理点的切空间，并解释 \(\partial(x^2)/\partial x=0\) 为什么不使该曲线具有二维切空间。
\end{exercise}
\begin{solution}
点为 \((a,a^2)\)，Jacobian 行为 \((0,1)\)，所以切空间是 \(\bigl\{(v,0):v\in k\bigr\}\)，维数一。方程对 \(y\) 的偏导仍为单位，给出非零约束。事实上 \(B\cong k[x]\)，微分关系只是 \(\mathrm{d}y=0\)，余下 \(\mathrm{d}x\) 自由；不能只检查某一个偏导数。
\end{solution}

\begin{exercise}\label{b3:ex:14-three-planes-tangent}
设 \(k\) 为域，\(B=k[x,y,z]/(xyz)\)。求每个 \(k\) 有理点的切空间维数，按恰有一个坐标为零和至少两个坐标为零分类。
\end{exercise}
\begin{solution}
Jacobian 行为 \((yz,xz,xy)\)。若恰有一个坐标为零，另外两个非零，三项中恰有一项非零，矩阵秩一，切空间维数二；若至少两个坐标为零，整行全零，切空间维数三。满足方程的点至少有一个零坐标，这两类穷尽所有有理点。计算不需要对 \(k\) 的特征加限制。
\end{solution}

\begin{exercise}\label{b3:ex:14-fat-point-tangent}
设 \(k\) 为域，\(r\ge2\)，\(B=k[x,y]/(x,y)^r\)。求唯一有理点处的切空间维数及 \(\dim_kB\)。说明前者能否恢复后者。
\end{exercise}
\begin{solution}
唯一极大理想为 \(\mathfrak m=(x,y)\)，而 \(\mathfrak m/\mathfrak m^2\) 的基为 \(\bar x,\bar y\)，所以切空间维数二。\(B\) 的基由全次数小于 \(r\) 的单项式组成，故
\(\dim_kB=\sum_{d=0}^{r-1}(d+1)=r(r+1)/2\)。不同 \(r\) 给出相同切空间维数而不同的代数维数，说明切空间只保存一阶信息。
\end{solution}

\begin{exercise}\label{b3:ex:14-cusp-tangent-map}
设 \(k\) 为域，考虑尖点参数化 \(\nu:t\mapsto(t^2,t^3)\)。直接求 \(t=a\) 处诱导的切空间映射。证明它在 \(a=0\) 处为零，在 \(a\ne0\) 处单射。
\end{exercise}
\begin{solution}
把 \(a+v\varepsilon\) 代入得到
\((a^2+2av\varepsilon,a^3+3a^2v\varepsilon)\)，故切映射为
\(v\mapsto(2av,3a^2v)\)。它满足尖点的切方程，因为
\(-3a^4(2av)+2a^3(3a^2v)=0\)。在零点两坐标均为零；若 \(a\ne0\)，整数二和三不可能同时在域中为零，至少一项系数非零，故映射单射。参数化在点集上单射不保证其每一点的切映射单射。
\end{solution}

\begin{exercise}\label{b3:ex:14-characteristic-two-circle}
设 \(k\) 为域。比较 \(B=k[x,y]/(x^2+y^2-1)\) 在特征不为二和特征二时的微分模及有理点切空间。在特征二时求该环的约化商。
\end{exercise}
\begin{solution}
总有 \(\Omega_{B/k}=(B\,\mathrm{d}x\oplus B\,\mathrm{d}y)/(2x\cdot\mathrm{d}x+2y\cdot\mathrm{d}y)\)。特征不为二时，在每个有理点 \((a,b)\) 上，\(a,b\) 不同时为零，故切空间维数一。特征二时，关系多项式为 \((x+y-1)^2\)，微分关系为零，所以微分模自由秩二，所有有理点的切空间维数二。约化商是 \(k[x,y]/(x+y-1)\cong k[x]\)，它在对应点的切空间维数一。这说明改变特征可能把一个方程变成重方程。
\end{solution}

\begin{exercise}\label{b3:ex:14-separable-derivation-extension}
设 \(k\subseteq K\)、\(L=K(\alpha)\)，其中 \(\alpha\) 对 \(K\) 可分，最小多项式为 \(f(T)=\sum_i a_iT^i\)。给定 \(k\) 导子 \(D:K\to L\)，证明它唯一延拓为 \(\widetilde D:L\to L\)，并写出 \(\widetilde D(\alpha)\)。
\end{exercise}
\begin{solution}
求导 \(f(\alpha)=0\) 迫使
\[
\widetilde D(\alpha)=-f'(\alpha)^{-1}\sum_iD(a_i)\alpha^i.
\]
右边有意义，因为可分性使 \(f'(\alpha)\ne0\)。在 \(K[T]\) 上，把系数按 \(D\) 求导、把 \(T\) 的导数指定为该值，得到取值于 \(L\) 的导子，它零化 \(f\)。对任意 \(hf\)，Leibniz 法则使导数在代入 \(\alpha\) 后也为零，所以它下降到 \(K[T]/(f)=L\)。其在系数及生成元上的值均已被迫确定，故唯一。
\end{solution}

\begin{exercise}\label{b3:ex:14-purely-inseparable-tower}
设 \(p\) 为素数，\(s\) 为超越元，令 \(k=\mathbb F_p(s)\)，\(L=k(\alpha)\)、\(E=k(\beta)\)，其中 \(\alpha^p=s\)、\(\beta^p=\alpha\)。计算复合微分正合列的三个非零候选项和两条箭头。
\end{exercise}
\begin{solution}
有 \(\Omega_{L/k}=L\,\mathrm{d}\alpha\)、\(\Omega_{E/k}=E\,\mathrm{d}\beta\)、\(\Omega_{E/L}=E\,\mathrm{d}\beta\)。第一箭头将 \(1\otimes \mathrm{d}\alpha\) 送到 \(\mathrm{d}(\beta^p)=0\)，故为零；第二箭头将 \(\mathrm{d}\beta\) 送到 \(\mathrm{d}\beta\)，为同构。因此该列为
\(E\xrightarrow{0}E\xrightarrow{1}E\to0\)。即使各域扩张次数都有限，首箭头仍可完全丢失一个微分方向。
\end{solution}

\begin{exercise}\label{b3:ex:14-perfect-closure}
设 \(p\) 为素数，\(t_1,t_2,\ldots\) 在 \(\mathbb F_p\) 上代数独立，令
\[
L=\mathbb F_p(t_1,t_2,\ldots),\qquad
K=\mathbb F_p(t_1^p,t_2^p,\ldots).
\]
证明 \(L/K\) 是无限次、指数为一的纯不可分扩张，并证明 \(\Omega_{L/K}\) 是以 \(\mathrm{d}t_i\)、\(i\ge1\)，为基的自由 \(L\) 模。比较此例与正文无限根塔的微分消失现象，说明无限次纯不可分扩张的微分模不必为零。
\end{exercise}
\begin{solution}
每个有理函数只涉及有限多个 \(t_i\)，其 \(p\) 次幂属于 \(K\)，所以扩张纯不可分、指数至多一；\(t_1\notin K\)，故指数恰为一。对任意 \(r\ge1\)，单项式 \(t_1^{e_1}\cdots t_r^{e_r}\)、\(0\le e_i<p\)，对 \(K\) 线性无关：一个假定关系只涉及有限多个变量和分母，清分母后，各项在前 \(r\) 个变量上的指数模 \(p\) 属于不同余数类，不会相消。它们也生成 \(K(t_1,\ldots,t_r)\)，故这个子扩张的次数为 \(p^r\)，从而 \([L:K]\) 无限。

所有 \(\mathrm{d}t_i\) 生成 \(\Omega_{L/K}\)，因为每个元素的微分都能按分式法则写成有限组合。对固定 \(j\)，形式偏导 \(\partial/\partial t_j\) 唯一延拓到 \(L\)，并零化 \(K\)。将它对应的线性映射作用于任何有限关系 \(\sum_i c_i\cdot\mathrm{d}t_i=0\)，便得 \(c_j=0\)。因此
\[
\Omega_{L/K}\cong\bigoplus_{i\ge1}L\,\mathrm{d}t_i.
\]
正文根塔使每个元素都在下一层成为 \(p\) 次幂；本题则不断加入新的独立方向，却不再对它们取 \(p\) 次根，两种无限扩张保留的微分因而不同。
\end{solution}

\begin{exercise}\label{b3:ex:14-arithmetic-etale}
令 \(B=\mathbb Z[t]/(t^2-t-1)\)。证明 \(B[1/5]\) 在 \(\mathbb Z[1/5]\) 上平展，并证明特征五纤维不是非分歧代数。
\end{exercise}
\begin{solution}
首一除法给出基 \(1,t\)，故自由平坦。微分关系为 \((2t-1)\cdot\mathrm{d}t=0\)，而在 \(B\) 中
\((2t-1)^2=4t^2-4t+1=5\)。倒置五后 \(2t-1\) 为单位，所以微分模为零，有限展示条件也成立。模五后，\(t^2-t-1=(t-3)^2\)，纤维为 \(\mathbb F_5[\varepsilon]/(\varepsilon^2)\)，其微分模为 \(\bigl(\mathbb F_5[\varepsilon]/(\varepsilon)\bigr)\mathrm{d}\varepsilon\ne0\)，故不非分歧。
\end{solution}

\begin{exercise}\label{b3:ex:14-square-zero-newton}
在 \(C=\mathbb Q[\varepsilon]/(\varepsilon^2)\) 中，求满足 \(x^2=1+\varepsilon\) 且 \(x\equiv1\pmod\varepsilon\) 的全部元素。若把 \(\mathbb Q\) 改成特征二的域，会发生什么？
\end{exercise}
\begin{solution}
写 \(x=1+a\varepsilon\)，则 \(x^2=1+2a\varepsilon\)，唯一解为 \(1+\varepsilon/2\)。这正是对误差 \(-\varepsilon\) 用导数二作一次修正。特征二时，所有这样的平方都等于一，因此没有解；不是有多个解，而是给定提升根本不存在。
\end{solution}

\begin{exercise}\label{b3:ex:14-free-differentials-not-smooth}
设 \(k\) 为特征 \(p>0\) 的域，\(A=k[s]\)、\(B=A[u]/(u^p-s)\)。证明 \(B\) 是自由秩 \(p\) 的 \(A\) 模，\(\Omega_{B/A}\) 是自由秩一的 \(B\) 模，但 \(A\to B\) 不形式光滑。说明这比“微分模自由但不平坦”的反例多检验了哪个条件。
\end{exercise}
\begin{solution}
首一除法给出 \(A\) 基 \(1,u,\ldots,u^{p-1}\)，所以 \(B\) 对 \(A\) 平坦。相对微分中 \(\mathrm{d}s=0\)，而 \(\mathrm{d}(u^p)=0\)，故关系不增加微分约束，得到 \(\Omega_{B/A}=B\,\mathrm{d}u\)。

置 \(f=T^p-s\)、\(C=A[T]/(f^2)\)、\(J=(f)/(f^2)\)，则 \(J^2=0\)、\(C/J=B\)。若恒等同态 \(B\to C/J\) 有 \(A\) 代数提升，\(u\) 的像必为 \(T+h\)，其中 \(h\in J\)。但
\[
(T+h)^p-s=f+h^p=f\ne0\quad\text{于}\;C,
\]
因为 \(p\ge2\) 使 \(h^p=0\)，而 \(f\notin(f^2)\)。故提升不存在。这里有限展示、平坦性与微分模自由都已成立，仍不能替代形式光滑所要求的提升存在性。
\end{solution}

\begin{exercise}\label{b3:ex:14-smooth-hyperbola-lifting}
设 \(A\) 为交换环，\(B=A[x,y]/(xy-1)\)。直接证明 \(A\to B\) 形式光滑；并在一个明确的平方零扩张中构造同一映射的两个不同提升，说明它一般不形式非分歧。
\end{exercise}
\begin{solution}
给定 \(B\to C/J\)，其中 \(J^2=0\)，\(x\) 的像为单位。任取它的提升 \(u\in C\)，仍为单位，令 \(x\mapsto u\)、\(y\mapsto u^{-1}\)，即得提升。为展示不唯一，取 \(A=k\) 为域，\(C=k[\varepsilon]/(\varepsilon^2)\)，基点为 \(x,y\mapsto1\)。一组提升为 \((1,1)\)，另一组为 \((1+\varepsilon,1-\varepsilon)\)，乘积均为一，且两组不同。
\end{solution}

\begin{exercise}\label{b3:ex:14-formal-base-change}
设交换环同态 \(A\to B\) 光滑、非分歧或平展。对任意交换 \(A\) 代数 \(A'\)，记 \(B'=A'\otimes_AB\)。利用\secref{b3:sec:1405}的形式提升性质，分别验证 \(A'\to B'\) 的同一性质；写出有限型、有限展示及平坦性在验证中的作用。
\end{exercise}
\begin{solution}
若 \(B=A[b_1,\ldots,b_n]\)，则 \(B'\) 由 \(1\otimes b_i\) 作为 \(A'\) 代数生成，故有限型保持。若 \(B=A[X_1,\ldots,X_n]/(f_1,\ldots,f_s)\)，则
\[
B'\cong A'[X_1,\ldots,X_n]/(f'_1,\ldots,f'_s),
\]
其中 \(f'_i\) 由结构同态 \(A\to A'\) 变换系数，故有限展示保持。光滑情形再用正文已证的形式光滑基变换性质。非分歧情形用有限型及
\(\Omega_{B'/A'}\cong A'\otimes_A\Omega_{B/A}=0\)。平展情形还需有限展示以及\cref{b3:prop:flat-base-change-transitivity}给出的平坦性保持。三种性质所需的条件都分别保留下来。
\end{solution}

\begin{exercise}\label{b3:ex:14-formal-composition}
设交换环同态 \(A\to B\to D\) 中的两个映射均形式光滑，或均形式非分歧，或均形式平展。证明复合映射分别具有同一性质。
\end{exercise}
\begin{solution}
给定 \(A\) 代数平方零提升问题 \(D\to C/J\)，先限制到 \(B\)。形式光滑时先提升 \(B\to C\)，把 \(C\) 视为这个结构下的 \(B\) 代数，再利用 \(B\to D\) 的形式光滑性提升整个 \(D\) 映射。形式非分歧时，两个 \(D\) 提升限制到 \(B\) 必相同，因此是同一 \(B\) 代数结构下的两个提升，再由第二映射的唯一性相等。两种性质同时成立便得到形式平展的结论。
\end{solution}
